% !TEX root = main.tex

\section{Results and Discussion}
\label{sec:results}

All results are computed on the same held-out block of 6\,008 observations
covering 2023--2025, against the same per-ticker mean benchmark, and every
comparison is tested rather than asserted.

\subsection{The conditional mean}
\label{sec:mean}

No model class improves on the per-ticker mean benchmark. Under the
fixed-origin protocol, pooled ridge attains an out-of-sample $R^2$ of
$-0.11\%$ and the pooled random forest $-0.34\%$; under the expanding-window
protocol the corresponding figures are $-0.28\%$ and $-0.31\%$. Diebold--Mariano
tests against the benchmark return $p$-values of $0.36$, $0.18$, $0.17$ and
$0.25$ respectively. None of the four is distinguishable from a forecast that
simply reports each ticker's historical average.

Directional accuracy tells the same story from a different angle. The benchmark
attains $0.5426$, which is exactly the proportion of positive days in the test
window, since a constant positive forecast is right whenever the market rises.
No model exceeds it.

Table~\ref{tab:per-ticker} disaggregates by ticker, and this is where the
evidence is strongest. Across 48 model--ticker Diebold--Mariano tests, exactly
one rejects at the uncorrected 5\% level: the hierarchical model on NVDA, at
$p=0.0216$. That single rejection is fewer than the 2.4 that 48 independent
tests would produce by chance under the null, and it survives neither
Benjamini--Hochberg control of the false discovery rate
\citep{benjamini1995controlling} nor Bonferroni correction. Under the
expanding-window protocol there are no rejections at all, corrected or
otherwise. The pattern is not a failure to detect predictability; it is what a
sample containing none looks like.

\subsection{Sequence models and the value of regularisation}
\label{sec:neural}

A natural objection is that linear and tree models are too rigid to find a
nonlinear signal. Four sequence architectures---LSTM, GRU, a one-dimensional
convolutional network and a Transformer encoder---were therefore estimated on
the same features, with the same training budget, and with the same privilege
ridge receives: each searched six regularisation configurations and selected one
by validation loss.

The outcome is unanimous. All four select the strongest available weight decay
and converge on a forecast whose standard deviation is zero to three decimal
places. Out-of-sample $R^2$ ranges from $-0.04\%$ to $-0.09\%$, and
Diebold--Mariano $p$-values against the benchmark range from $0.13$ to $0.51$.
Given free choice of regularisation on held-out data, every architecture
concludes that the optimal amount of signal to extract from fourteen
macro-financial predictors is none, and collapses to the benchmark.

Two caveats belong with this result. The selected penalty sits at the top of the
search grid, which would ordinarily indicate too narrow a search; here it does
not, because the forecast standard deviation has already reached zero and no
larger penalty can alter it. And the finding is conditional on this information
set, this decade and these eight assets. It is not a claim about market
efficiency in general, but a demonstration consistent with the local and
episodic predictability described by \citet{timmermann2008elusive}.

The contrast with the untuned specifications is instructive in its own right.
Without a regularisation search the same four architectures return out-of-sample
$R^2$ between $-6.5\%$ and $-35.6\%$, with forecast dispersion between
$23\%$ and $55\%$ of the realised return standard deviation. The unregularised
models are not finding signal; they are fitting noise and paying for it out of
sample.

\subsection{The conditional variance}
\label{sec:variance}

The second moment behaves entirely differently. Table~\ref{tab:volatility}
reports QLIKE for the three conditional specifications and a 20-day rolling
variance benchmark, all scored on an identical sample.

All three conditional models beat the benchmark decisively: GARCH(1,1)-$t$ at
$\mathrm{DM}=-4.33$ ($p=0.00002$), GJR at $-3.48$ ($p=0.0005$) and EGARCH at
$-3.32$ ($p=0.0009$). This is the first comparison in the study in which a
richer model reliably outperforms a naive alternative.

Among the conditional specifications, however, the differences are largely not
significant. GARCH attains the best QLIKE at $-6.536$, but against EGARCH the
difference carries $p=0.325$ and against GJR it sits on the boundary at
$p=0.0499$; GJR against EGARCH gives $p=0.687$. The asymmetric extensions,
motivated by the leverage effect \citep{nelson1991conditional,
glosten1993relation}, deliver no measurable improvement for these tickers over
this window. The symmetric specification is retained on parsimony grounds.

\subsection{Decomposing the probabilistic gain}
\label{sec:decomposition}

Sections~\ref{sec:mean} and~\ref{sec:variance} point in opposite directions, and
Figure~\ref{fig:decomposition} resolves them. A predictive distribution has a
location, a scale and a shape; the ladder adds one at a time to a null
predictive, scoring each rung on the same sample.

\begin{center}
\begin{tabular}{lrrr}
\toprule
Addition & Gain (nats) & Share & DM $p$ \\
\midrule
Conditional variance (GARCH)      & $+0.1528$ & 54.4\% & $<0.00001$ \\
Conditional mean (ridge)          & $+0.0020$ &  0.7\% & $0.0293$ \\
Fat tails and scale calibration   & $+0.1261$ & 44.9\% & $<0.00001$ \\
\midrule
Total                             & $+0.2808$ & 100\%  & \\
\bottomrule
\end{tabular}
\end{center}

The conditional mean contributes $0.7\%$ of the total gain. That contribution is
statistically detectable at $p=0.029$ and economically trivial, and the two
statements are not in tension: with 6\,008 observations a Diebold--Mariano test
resolves differences far below any that would affect a decision. The distinction
is worth stating explicitly, because a study reporting only the $p$-value would
present this as evidence that macroeconomic conditioning improves density
forecasts, which it does not in any practical sense.

The remaining $99.3\%$ divides between the conditional variance ($54.4\%$) and
the shape of the predictive distribution ($44.9\%$). Getting the shape right is
worth nearly as much as getting the scale right---a result that would be
invisible to any evaluation confined to point forecasts.

\subsection{Calibration}
\label{sec:calibration-results}

Table~\ref{tab:probabilistic} and Figures~\ref{fig:reliability}
and~\ref{fig:pit} assess whether the stated uncertainty is honest.

The Gaussian predictive is well calibrated in the tails, attaining $0.909$
against a $90\%$ nominal level and $0.943$ against $95\%$, but over-covers
substantially in the body at $0.572$ against $50\%$. This is the signature of a
thin-tailed density fitted to fat-tailed returns: with no allowance for
occasional large moves, the density widens everywhere to accommodate them.

Student-$t$ innovations with validation-based scale calibration attain $0.536$,
$0.825$, $0.918$ and $0.958$ against the four nominal levels, a mean absolute
coverage error of $0.022$ against the Gaussian's $0.030$, alongside a decisively
better log score ($2.478$ against $2.352$, $p<0.00001$) and CRPS.

Two further points deserve emphasis, both of which qualify the result.

First, the Gaussian specification concealed a scale error rather than avoiding
it. Measured as the standard deviation of $y/\sigma$, the scale the data require
is $1.002$ on the training block but $1.117$ on the test block: GARCH parameters
frozen at end-2019 understate 2023--2025 volatility by roughly twelve percent.
The Gaussian likelihood inflates its variance to absorb outliers, and that
inflation happened to offset the shortfall, producing accidentally reasonable
$90\%$ coverage alongside visibly over-wide $50\%$ intervals. Separating tail
behaviour from scale makes the miscalibration visible, and the validation block
corrects it: the scale ratio there is $1.118$, almost exactly the test value.

Second, calibration is improved rather than achieved. The Kolmogorov--Smirnov
statistic against PIT uniformity falls from $0.0517$ under the Gaussian to
$0.0206$ under the Student-$t$, but the latter still rejects uniformity at
$p=0.012$. The predictive distribution is materially better, not correct, and
Section~\ref{sec:conclusion} treats the residual misspecification as open.

\subsection{What survives shrinkage}
\label{sec:selection}

Figure~\ref{fig:shrinkage} reports the horseshoe shrinkage weight for each
predictor. One of fourteen has a $95\%$ credible interval excluding zero: the
daily change in the effective federal funds rate, at
$0.00063$ with interval $[0.00028,\,0.00097]$ and shrinkage weight $0.89$. The
market return and the change in the VIX follow at weights $0.67$ and $0.56$,
with intervals covering zero. Every own-return lag is shrunk toward zero.

The direction is consistent with the monetary-policy transmission documented by
\citet{bernanke2005explains}, and the magnitude is economically small: roughly
six basis points of next-day return per standard-deviation change in the policy
rate, against a daily return standard deviation of 248 basis points.

The contrast with the v1 correlation screen is the substantive point. That
procedure selected the ten predictors most correlated with the target on the
training block, and a trending series correlates with almost any target over a
four-year window; it therefore preferentially selected the non-stationary levels
that Section~\ref{sec:extrapolation} shows to be destructive. Placing selection
inside the posterior removes that failure mode.

\subsection{Cross-asset transfer}
\label{sec:transfer}

Table~\ref{tab:transfer} reports leave-one-ticker-out performance. A ticker
entirely absent from estimation, whose intercept is integrated over the
population distribution, attains a mean log score of $2.4727$ against $2.4777$
for the same model estimated with that ticker included---a gap of $0.005$
nats---and retains $87\%$ coverage at the $90\%$ nominal level.

Revealing the held-out ticker's own history changes almost nothing. Across
$k \in \{0, 20, 60, 120, 250\}$ days the mean log score moves from $2.4727$ to
$2.4722$ and RMSE is unchanged to five decimal places, while the posterior
demonstrably shifts between settings.

The reading is that these eight assets are close to exchangeable for forecasting
purposes: the between-ticker intercept variance is small relative to observation
noise, so the population layer supplies essentially everything and asset
identity supplies almost nothing. This is what makes partial pooling worth its
complexity here, since it delivers usable predictive intervals for an asset with
no estimation history at all.

\subsection{Fixed origin against expanding window}
\label{sec:extrapolation}

The final result is methodological, and it is reported because the failure it
describes is silent.

An earlier version of this pipeline selected model inputs by column type, which
admitted the three non-stationary levels of Section~\ref{sec:features}.
Standardised on 2016--2019 moments and applied to 2023--2025, the S\&P~500 level
arrives at $z = +8.7$, with $100\%$ of test observations beyond three training
standard deviations; the policy rate arrives at $z=+4.7$ and the ten-year yield
at $z=+3.9$ (Figure~\ref{fig:extrapolation}a).

Every model with an unbounded output head extrapolated along those axes. The
LSTM's test predictions correlate with the S\&P~500 level at $r=0.88$ and imply a
$+12\%$ return every trading day for three years. Sixteen of twenty-one saved
models failed a prediction sanity check. Training error remained normal
throughout, no exception was raised, and metric files were written as usual.

Replacing the levels with stationary counterparts moves mean out-of-sample $R^2$
across the eight tickers from $-6.23\%$ to $-0.19\%$ for ridge, from $-26.16\%$
to $-0.47\%$ for the random forest and from $-24.58\%$ to $+0.13\%$ for the
hierarchical model (Figure~\ref{fig:extrapolation}b). For the sequence models
the corresponding shift is from per-ticker means between $-638\%$ and $-10\,425\%$
to between $-0.04\%$ and $-0.09\%$.

Two general lessons follow. Fitting a scaler once and freezing it is safe only
for stationary predictors, and the expanding-window protocol, which refits the
scaler inside each fold, is the appropriate default when that cannot be
guaranteed: it alone moves fixed-split ridge from $-0.044$ to $-0.013$ in
out-of-sample $R^2$ on the uncorrected feature set. More importantly, the
failure produces no diagnostic of its own. It is detectable only by auditing the
distribution of predictors across splits and the scale of the predictions
themselves, which is why both audits are released with the code.
